\chapter{CONCLUSIONES}
%\section{Conclusiones}

Se concluye que la implementación de un modelo predictivo de demanda contribuye a mejorar la toma de decisiones en el proceso de pedido a proveedores en Donna S.A.C., debido a que permite estimar la demanda futura de productos perecibles a partir de datos históricos de ventas y variables temporales relevantes. De esta manera, la empresa puede contar con una herramienta de apoyo para planificar el reabastecimiento con mayor precisión, reduciendo la dependencia de decisiones basadas únicamente en la experiencia o criterios subjetivos.

Respecto al primer objetivo específico, se determinó que la calidad de la base de datos histórica influye directamente en la precisión de los modelos predictivos. La disponibilidad de registros organizados, completos y consistentes permitió construir una serie temporal adecuada para el análisis de la demanda. Asimismo, la identificación de valores atípicos, datos duplicados o inconsistencias fue necesaria para evitar distorsiones en el entrenamiento del modelo y mejorar la confiabilidad de los resultados obtenidos.

En relación con el segundo objetivo específico, se concluye que las técnicas de preprocesamiento de datos fueron fundamentales para mejorar el rendimiento de los modelos de predicción. La limpieza de datos, la agrupación diaria de ventas, el tratamiento de valores atípicos, la selección de columnas relevantes y la construcción de variables temporales permitieron transformar los registros históricos en una base adecuada para el modelamiento predictivo. Estas actividades fortalecieron la capacidad del modelo para identificar patrones de comportamiento en la demanda de productos perecibles.

Con respecto al tercer objetivo específico, se identificó que las variables temporales e históricas tienen una influencia relevante en el comportamiento de la demanda. Factores como el día de la semana, el mes, las ventas históricas, los rezagos de demanda y los promedios móviles permitieron representar mejor las variaciones de consumo. Esto confirma que la demanda de productos perecibles no presenta un comportamiento uniforme, sino que responde a patrones temporales que deben ser considerados en el proceso de pronóstico.

En cuanto al cuarto objetivo específico, se evaluaron distintos modelos de forecasting y machine learning mediante métricas como MAE, MSE, RMSE, MAPE y R$^2$. A partir de esta comparación, se seleccionó el modelo con mejor desempeño predictivo y mayor aplicabilidad para el contexto de la empresa. Según los resultados obtenidos, el modelo seleccionado presentó una capacidad adecuada para aproximarse a la demanda real, por lo que puede ser utilizado como apoyo para la planificación de compras y la optimización del reabastecimiento.

Finalmente, se concluye que la propuesta desarrollada representa una mejora metodológica y operativa para Donna S.A.C., ya que permite convertir los datos históricos de ventas en información útil para la toma de decisiones. La implementación del modelo predictivo no solo aporta a la estimación de la demanda futura, sino que también contribuye a reducir riesgos asociados al sobrestock, quiebres de stock y pérdidas por productos perecibles, fortaleciendo así la eficiencia del proceso de abastecimiento.




%%%%%CAMBIO
% Se puede concluir que el desarrollo de modelos predictivos de demanda, utilizando técnicas avanzadas como Regresión Lineal Multiple, es esencial para mejorar la rentabilidad en la venta de productos perecibles. Integrar datos históricos de ventas con variables clave como feriados, pago proveedores y precios no solo facilita una gestión de inventario más precisa, sino que también optimiza las estrategias de precios y promociones, lo que resulta en una planificación más efectiva y en un aumento de la rentabilidad.

% La investigación enfocada en modelos de Forecasting para maximizar la planificación de recursos en Donna S.A.C. destaca la importancia de recopilar datos detallados sobre el desempeño de ventas en diversas áreas y períodos. Estos datos son cruciales para identificar patrones de comportamiento de los clientes, evaluar el impacto de estrategias comerciales variadas y anticipar cambios en la demanda. La implementación adecuada de modelos de Forecasting no solo facilita la predicción precisa de ventas futuras, sino que también proporciona valiosos insights para la toma de decisiones estratégicas que impulsan el crecimiento sostenido de la empresa.

% El análisis de modelos de Forecasting, al considerar la estructura temporal de los datos y factores como tendencias, estacionalidad y variación aleatoria, permite la aplicación de modelos estadísticos avanzados para pronosticar con precisión el comportamiento futuro de las ventas. Este enfoque es fundamental para empresas como Donna S.A.C., donde la capacidad de prever la demanda de manera confiable puede marcar la diferencia en la planificación operativa y estratégica. Así, el uso efectivo de modelos de Forecasting no solo mejora la eficiencia en la gestión de recursos, sino que también fortalece la capacidad de la empresa para adaptarse ágilmente a cambios en el mercado y maximizar su rentabilidad.


%
%
%
%\section{Recomendaciones}
%
%Nisi porta lorem mollis aliquam ut porttitor leo. Aenean pharetra magna ac placerat vestibulum. Est placerat in egestas erat imperdiet sed euismod. Velit euismod in pellentesque massa placerat. Enim praesent elementum facilisis leo vel fringilla. Ante in nibh mauris cursus mattis molestie a iaculis. Erat pellentesque adipiscing commodo elit at imperdiet dui accumsan sit. Porttitor lacus luctus accumsan tortor posuere ac ut. Tortor at auctor urna nunc id. A iaculis at erat pellentesque adipiscing commodo elit. 


